\begin{ouverture}[grandeurs]{14}{Durées : jusqu’à la seconde}{Partie 3 · Grandeurs et mesures}
{Heures · Minutes · Secondes · Calculs d’instants}
\objectif{Convertir entre heures, minutes et secondes.}
\objectif{Calculer une durée, un début ou une fin d’évènement.}
\objectif{Résoudre des problèmes de temps en plusieurs étapes.}
\prerequis{Durées en heures et minutes du CM1.}
\end{ouverture}
\pagestyle{otMain}\essentiel\fontsize{11}{15}\selectfont
\phantomsection\label{cours-14-1}
\essSection{1}{La seconde entre en scène}
\begin{definition}Une minute vaut 60 secondes et une heure vaut 60 minutes, donc une heure vaut $60\times60=3\,600$ secondes. Les durées s’écrivent en heures, minutes et secondes : 1 h 24 min 35 s.\end{definition}
\begin{savoirfaire}[Des paquets de 60]Convertir : 2 min 15 s $=120+15=135$ s ; 200 s $=180+20=3$ min 20 s. Pour les heures : 1 h 12 min $=72$ min ; 100 min $=1$ h 40 min. On échange toujours des paquets de 60.\end{savoirfaire}
\begin{exemple}Le chronomètre affiche souvent les secondes : au 100 mètres, 58 s contre 1 min 01 s, c’est 58 contre 61 : trois secondes d’écart. En natation, 1 min 24 s 5 dixièmes se lit aussi 84,5 secondes.\end{exemple}
\begin{attention}Un horaire est un instant, une durée est un temps écoulé : 60 minutes font une heure, jamais « 100 minutes » dans une heure. Le système du temps est en base 60, pas en base 10.\end{attention}
\clearpage
\phantomsection\label{cours-14-2}
\essSection{2}{Calculer avec les durées}
\begin{savoirfaire}[Avancer par étapes rondes]Durée entre deux instants : avancer par étapes. De 10 h 47 min 30 s à 11 h 02 min 10 s : 12 min 30 s pour atteindre 11 h, puis 2 min 10 s : total 14 min 40 s.\end{savoirfaire}
\begin{exemple}Additionner des durées : 2 min 40 s $+$ 1 min 35 s $=3$ min 75 s $=4$ min 15 s : les 75 secondes s’échangent contre 1 min 15 s. La retenue vaut 60, pas 100.\end{exemple}
\begin{propriete}Heure de fin $=$ heure de début $+$ durée ; heure de début $=$ heure de fin $-$ durée. Une frise des instants évite les erreurs de sens.\end{propriete}
\begin{attention}Comparer des durées exige la même unité : 1 min 30 s et 95 s se comparent en secondes ($90<95$). Les dixièmes de seconde départagent les courses serrées.\end{attention}

\clearpage\rubrique{Découverte}
\begin{blocLivre}{ex-14-1}\begin{exercice}[1]{Convertir en secondes}
Convertir : 2 min ; 2 min 15 s ; 5 min 08 s.
\end{exercice}\end{blocLivre}
\begin{blocLivre}{ex-14-2}\begin{exercice}[1]{Convertir en minutes et secondes}
Convertir en minutes et secondes : 200 s ; 428 s ; 59 s ; 60 s.
\end{exercice}\end{blocLivre}
\begin{blocLivre}{ex-14-3}\begin{exercice}[1]{L’heure en secondes}
Combien de secondes dans une heure ? Dans un quart d’heure ?
\end{exercice}\end{blocLivre}
\begin{blocLivre}{ex-14-4}\begin{exercice}[1]{Additionner des durées}
Calculer : 2 min 40 s $+$ 1 min 35 s, puis 3 h 45 min $+$ 2 h 30 min.
\end{exercice}\end{blocLivre}
\begin{blocLivre}{ex-14-5}\begin{exercice}[1]{Le chrono du 100 m}
Aux essais, Lina court le 100 m en 1 min 01 s et Noa en 58 s. Qui est le plus rapide, et de combien ?
\end{exercice}\end{blocLivre}

\rubrique{Entraînement}
\begin{blocLivre}{ex-14-6}\begin{exercice}[2]{La durée du trajet}
Un métro part à 10 h 47 min 30 s et arrive à 11 h 02 min 10 s. Combien dure le trajet ?
\end{exercice}\end{blocLivre}
\begin{blocLivre}{ex-14-7}\begin{exercice}[2]{L’heure de fin}
Un lavage de voiture commence à 9 h 52 et dure 25 min. Un cycle de machine commence à 17 h 48 min 30 s et dure 1 min 40 s. Donner les deux heures de fin.
\end{exercice}\end{blocLivre}
\begin{blocLivre}{ex-14-8}\begin{exercice}[2]{Remonter le temps}
Un gâteau doit sortir du four à 18 h 10 après 35 min de cuisson. À quelle heure faut-il l’enfourner ?
\end{exercice}\end{blocLivre}
\begin{blocLivre}{ex-14-9}\begin{exercice}[2]{Comparer des durées}
Quelle durée est la plus longue : 1 min 30 s ou 95 s ? Et 2 h 05 ou 118 min ?
\end{exercice}\end{blocLivre}
\begin{blocLivre}{ex-14-10}\begin{exercice}[2]{Les dixièmes qui décident}
En natation, Aya réalise $32{,}6$ s et Lou $32{,}48$ s. Qui touche le mur la première ? Quel est l’écart ?
\end{exercice}\end{blocLivre}

\rubrique{Problèmes}
\begin{blocLivre}{ex-14-11}\begin{exercice}[3]{La playlist}
Une chanson dure 3 min 45 s. Quelle est la durée de 4 écoutes ? Cette durée dépasse-t-elle un quart d’heure ?
\end{exercice}\end{blocLivre}
\begin{blocLivre}{ex-14-12}\begin{exercice}[3]{Le tour de circuit}
Un kart boucle un tour en 1 min 12 s. Combien de temps durent 5 tours ? Et combien de tours complets en 10 minutes ?
\end{exercice}\end{blocLivre}
\begin{blocLivre}{ex-14-13}\begin{exercice}[3]{La séance bien remplie}
Un cours de sport de 1 h 30 comprend 12 min d’échauffement, 55 min de jeu et des étirements pour finir. Combien de temps durent les étirements ?
\end{exercice}\end{blocLivre}
\begin{blocLivre}{ex-14-14}\begin{exercice}[3]{Le journal de bord}
Un randonneur marche de 9 h 40 à 12 h 05, s’arrête 45 min, puis marche de 12 h 50 à 15 h 10. Quelle est sa durée totale de marche ?
\end{exercice}\end{blocLivre}
\begin{blocLivre}{ex-14-15}\begin{exercice}[3]{Vrai ou faux ?}
Tom écrit : « 1 h 85 min, c’est une durée correcte » et « 2 min 30 s $=$ 2,3 min ». Corriger ces deux écritures.
\end{exercice}\end{blocLivre}

\clearpage
